\begin{textsample}{POS Dim 5 – global\_north\_2025\_04 – Score 15.00 – global\_north\_2025\_04/124064225.txt}  \label{ex:f5_pos_096}
Israel is starving Palestinians, envoy tells top UN \textbf{court} A Palestinian diplomat has told the United Nations ' top \textbf{court} that Israel is killing and displacing civilians and targeting aid workers in Gaza, in a \textbf{case} that Israel has criticised as part of its"systematic \textbf{persecution} and delegitimisation".
Israel denies deliberately targeting civilians and aid staff as part of its war with Hamas and did not attend the hearing at the \textbf{International} \textbf{Court} of Justice.
In The Hague, Palestinian Ambassador to the Netherlands Ammar Hijazi accused Israel of breaching \textbf{international} \textbf{law} in the occupied territories."Israel is starving, killing and displacing Palestinians while also targeting and blocking humanitarian organisations trying to save their lives,"he told the \textbf{court}.
The hearings are focused on a request last year from the UN General Assembly, which asked the \textbf{court} to weigh in on Israel ' s \textbf{legal} responsibilities after the country blocked the UN agency for Palestinian refugees from operating on its territory.
In a \textbf{resolution} sponsored by Norway, the General Assembly requested an \textbf{advisory} opinion @ @ @ @ @ @ @ @ @ @, on Israel ' s \textbf{obligations} in the occupied territories to"ensure and facilitate the \textbf{unhindered} provision of urgently needed supplies essential to the survival of the Palestinian civilian population".
The hearings opened as the humanitarian aid system in Gaza is nearing collapse.
Israel has blocked the entry of food, fuel, medicine and other humanitarian supplies since March 2.
It renewed its bombardment on March 18, breaking a ceasefire, and seized large parts of the territory, saying it aims to push Hamas to release more hostages.
Despite the stepped-up Israeli pressure, ceasefire efforts remain deadlocked.
The World Food Programme said last week that its food stocks in the Gaza Strip have run out under Israel ' s nearly eight-week-old blockade, ending a main source of sustenance for hundreds of thousands of Palestinians in the territory.
Many families are struggling to feed their children.
The United Nations was the first to address the \textbf{court} on Monday, followed by Palestinian representatives.
In total, 40 states and four \textbf{international} organisations are scheduled to participate. @ @ @ @ @ @ @ @ @ @ but could \textbf{submit} a written statement.
Israel ' s Foreign Ministry did not immediately respond to requests for comment.
The United States, which voted against the UN \textbf{resolution}, is scheduled to speak on Wednesday.
The \textbf{court} is likely to take months to deliver a \textbf{ruling}, but experts say the decision, though not legally binding, could profoundly affect \textbf{international} jurisprudence, \textbf{international} aid to Israel and public opinion."\textbf{Advisory} opinions provide clarity,"Juliette McIntyre, an expert on \textbf{international} \textbf{law} at the University of South Australia, told The Associated Press.
Governments rely on them in \textbf{international} negotiations and the outcome could be used to pressure Israel into easing restrictions on aid.
Whether any \textbf{ruling} will have an effect on Israel, however, is unclear.
Israel has long accused the United Nations of being unfairly biased against it and has ignored a 2004 \textbf{advisory} \textbf{ruling} by the ICJ that found its West Bank separation barrier illegal.
On Tuesday, South Africa, a staunch critic of Israel, will present its \textbf{arguments}.
In hearings @ @ @ @ @ @ @ @ @ @ the country accused Israel of committing \textbf{genocide} against the Palestinians in Gaza - a charge Israel denies.
Those \textbf{proceedings} are still under way.
Israel ' s ban on the UN relief agency, known as UNRWA, came into effect in January.
The organisation has faced increased criticism from Prime Minister Benjamin Netanyahu and his far-right allies, who claim the group is deeply infiltrated by Hamas.
UNRWA rejects that claim.
Israel alleged that 19 out of UNRWA ' s approximately 13,000 staff in Gaza took part in Hamas ' s attack in southern Israel on October 7 2023, which killed about 1,200 people and set off the war in Gaza.
UNRWA said it fired nine employees after an internal UN investigation concluded they could have been involved, although the evidence was not authenticated and corroborated.
Israel later alleged that about 100 other Palestinians in Gaza were Hamas members, but never provided any evidence to the United Nations.
Israel has also accused Hamas of using UN facilities for cover, building tunnels near UN buildings and diverting aid deliveries for its own @ @ @ @ @ @ @ @ @ @ Gaza, but it controls all entry to the territory, and its ban on UNRWA from operating inside Israel greatly limits the agency ' s ability to function.
Israeli officials say they are looking for alternative ways to deliver aid to Gaza that would cut out the United Nations.
UNRWA was established by the UN General Assembly in 1949 to provide relief for Palestinians who fled or were expelled from their homes in what is now Israel during the war surrounding Israel ' s creation the previous year until there is a political solution to the Israeli-Palestinian conflict.
The agency has been providing aid and services - including health and education - to some 2.5 million Palestinians in Gaza, the occupied West Bank and east Jerusalem, as well as three million more in Syria, Jordan and Lebanon.
Israel ' s air and ground war has killed more than 51,000 Palestinians, mostly women and children, according to the Gaza Health Ministry, which does not say how many of the dead were civilians or combatants.
Israel says it has killed around 20,000 militants @ @ @ @ @ @ @ @ @ @

% matched lemmas: advisory, argument, case, court, genocide, international, law, legal, obligation, persecution, proceeding, resolution, ruling, submit, unhindered
\end{textsample}
